% ============================================================================
% Title:        Materials compatibility
% Author: Nathaniel Thomas
% Affiliation:  Texas A&M University, Department of Nuclear Engineering
% Email:        nathaniel@tamu.edu
% Date Created: 3/15/2026
% Version:      1.0
%
% Description:
%   A .tex source file discussing materials compatability in molten salt
%   corrosion.
%
% =============================================================================

\section{Compatibility of Candidate Structural Materials}
\label{sec: compatability of candidate}

As stated previously, pure molten fluoride salts have very little to no
corrosive effects. This is due to the Gibbs energy of fluorination being
more favorable for the salt components (i.e. \ch{K}, \ch{Li}, \ch{Na},
and \ch{Be}) than for structural materials (i.e. \ch{Fe}, \ch{Cr},
\ch{Ni}, etc.). However, the presence of contaminants, either through
environmental contamination or fission products, will skew the redox
potential of the salt such that corrosion of structural materials
will occur. While the purification methods listed above can be used
to minimize the concentration of environmental contaminants (and even
fission products in the case of \ch{NF3}), they only work so well,
making selection of proper materials necessary. \par

\subsection{Ni Alloys}
\label{sec: ni alloys}

The simplest thing to do is select materials whose fluorination potential
is as far from the salt components as possible. To this end alloys made
of primarily \ch{Ni} or \ch{Mo} are practical, since they are the
materials that have a high Gibbs energy of fluorination and can be easily
welded. Since \ch{Ni} is the cheaper of the two metals it will be focused
on \cite{rev49}. While pure \ch{Ni} (Ni-200) would likely be the material
of choice from a pure chemistry standpoint, its poor structural strength
at process temperatures ($>\SI{500}{\celsius}$) makes its use limited.
The MSRE used pure \ch{Ni} liners in their purification vessels, and
although these corroded away when fluorine gas was used to purify used
salt, the concept is still sound \cite{rev21}. It simply requires
more active checking the liner between runs. Kelleher et. all used
a pressurized outer shell made of steel to add structural strength to
their Ni-200 purification vessel \cite{rev29}. \par

Ni alloy Hastelloy N (also known by INOR-8) is another appealing option
that has a lower but comparable corrosion resistance to Ni-200, due to
its high \ch{Ni} and \ch{Mo} content, with increased structural strength due
to small concentrations of \ch{W}, \ch{Fe}, and \ch{Cr} \cite{rev50,rev51}.
Hastelloy N was specifically made for the MSRE and was used as the sole
structural material in the MSRE loop \cite{rev51}. It showed almost no
corrosion across the entire run \cite{rev51}. Experimental data obtained by
Koger predicted the corrosion rate of Hastelloy N by dynamic FLiBe to be at
around \SI{0.02}{\milli\meter\per\year} \cite{rev36}. Experiments by Sandhi
et al. showed Hastelloy N to have a high resistance to corrosion by static,
molten FLiNaK salt \cite{rev37}. In their experiments, Hastelloy N
apparently suffered less surface corrosion, had a lower surface concentration
of \ch{Cr}, half the depth of surface penetration, and had a positive mass
change after spending \SI{100}{\hour} in \SI{700}{\celsius} FLiNaK, when
compared to stainless steel \cite{rev37}. Hastelloy N’s use in a purification
vessel is largely unproven, however its resistance to anhydrous \ch{HF} and
molten salts has been explored, and it proved to have a low rate of corrosion
compared to other high Ni alloys with a largely linearly changing corrosion
rate with temperature \cite{rev37,rev50,rev51}. \par

\subsection{Steels}
\label{sec: steels}

Steels due to their high structural strength at high temperatures and lower
cost should be at least considered. The main issue with most commercial
stainless steels is their high \ch{Cr} content, but \ch{Fe} is also more
easily fluorinated than \ch{Ni} or \ch{Mo}. If salts are sufficiently
purified then corrosion can be mitigated in a molten salt loop, although
as low of a chromium content as is possible is preferred. Steels however
cannot be used for creation of a purification vessel that utilizes \ch{HF}
or any stronger fluorinating agent as it will corrode too quickly.
Stainless steels can handle low temperature \ch{HF}, but they will not
handle it at process temperatures \cite{rev50}. For stronger fluorinating
agents like fluorine gas or higher temperature \ch{NF3}, all steels are
considered to have an unacceptable rate of corrosion. \cite{rev50} \par

\subsection{Copper}
\label{sec: copper}

A material that also historically has showed much promise is copper.
When scaling up to meet the high demands for purified salts, the
ARE swapped out their nickel reaction vessels for stainless steel vessels
lined with copper \cite{rev6}. This vessel design remained in intermediate
use with nickel lined vessels through to the conclusion of the MSRE
\cite{rev10,28}. They later made sampling vessels out of copper as well.
During the MSRE, when the nickel dip tubes corroded away, likely due to
the high sulfur content of received salt, attempted solutions involved
coating the nickel parts with copper \cite{rev28}. After the MSRE, the use of
copper in fluorination appears to have stopped, with studies into its
corrosion in molten salt environments only just recently resuming
\cite{rev60}. Despite this, a recent study by Weinstein et. all has shown
copper to have superior corrosion resistance with inferior chromium diffusion
resistance to nickel \cite{rev60}. Further studies into this may become
available soon. The main issue with this as a material is that, while copper
is highly resistant to dry fluorinating agents, it is readily corroded and
dissolved by hydrofluoric acid, so moisture must be purged readily to prevent
serious corrosion \cite{rev6}.

\subsection{Carbon}
\label{sec: carbon}

Although not mechanically useful, glassy carbon, graphite, and \ch{SiC}
do make for excellent crucible/liner materials in a reaction vessel,
since carbon is chemically inert to all but the harshest fluorinating
compounds (such as \ch{F2} and \ch{NF3}) \cite{rev31,rev61}. The main
issues of using carbon structures are that salt can diffuse into the
graphite \cite{rev62} and the increase in the corrosion of metals used
in other structures \cite{rev63,rev64}. Molten salt impregnation of
graphite is a well-documented phenomenon first observed in early ARE
materials tests \cite{rev61,rev65}. Depending on the pore size and void
fraction of the graphite, salt impregnation can cause serious structural
damage to the graphite via cracking, particularly during the frequent
cooling and heating cycles of a purifier \cite{rev62}. This makes glassy
carbon the preferred material to use over even high-quality graphite.
The increase in salt corrosivity due to graphite contamination is a
less understood phenomenon however it was discovered to occur in the MSRE
test loops \cite{rev63}. The exact transport mechanism of this phenomenon
is unknown and is currently under study; however, it is understood that
the carbon and metal particles do diffuse through the salt to form
carbide layers on both surfaces and that the formation of carbide
layers on metal surfaces is normally prevented by protective oxide layers
on the metals, which are destroyed in a molten salt or fluorinating
environment \cite{rev64}. These carbide layers can flake off and expose
fresh metal surfaces to be further carburized. \par
